\input{../tex-inputs/latex-leseansicht-vorspann}

\section[Gustav Schwarzkopf an Arthur Schnitzler, 10. 8. 1915]{L04245 Gustav Schwarzkopf an Arthur Schnitzler, 10. 8. 1915}
\nopagebreak\mylabel{L04245v}
\rehead{ }\normalsize\beginnumbering\briefempfaengerindex{Schnitzler, Arthur@\textsc{Schnitzler, Arthur}!zzzSchwarzkopf, Gustav@\emph{von Gustav Schwarzkopf}!1915-08-101@{10. 8. 1915}|(be}
\toendnotes[C]{\smallbreak\pagebreak[2]}
\correspDesc{Versand  durch Gustav Schwarzkopf am 10. 8. 1915 in Wien
\newline{}Erhalt  durch Arthur Schnitzler im Zeitraum [11. 8. 1915 – 15. 8. 1915?] in Bad Ischl}\toendnotes[C]{\smallbreak}
\Standort{CUL, Schnitzler, B 96a.}
\physDesc{Bildpostkarte, 570 Zeichen
\newline{}Handschrift: schwarze Tinte, deutsche Kurrent
\newline{}Versand: Stempel: »\nobreak{}\oindex{I., Innere Stadt@\textbf{I., Innere Stadt}, \emph{Verwaltungsgebiet}|pwk}1/\textsubscript{1} Wien 8, 10. VIII. 1\textcolor{gray}{5}, 3\nobreak{}«.  
\newline{}Schnitzler: mit Bleistift Vermerk: »\textsc{Schwarzkopf}« }\toendnotes[C]{\smallbreak}\pstart{}{\pb}Herrn\pend{}\pstart{}\textsc{D\textsuperscript{r} Arthur Schnitzler}\pend{}\pstart{}\textsc{Ischl (O. Ö.)\oindex{Bad Ischl@\textbf{Bad Ischl}|pw}}\pend{}\pstart{}\textsc{Hôtel Kaiserkrone\oindex{Hotel Kaiserkrone@\textbf{Hotel Kaiserkrone}, \emph{Hotel}|pw}.}\pend{}{\bigskip}
\pstart
           \noindent{}\centering{}{\pb}\textcolor{gray}{\textbf{Wien, 1. Schottenkirche}}\oindex{Wien@\textbf{Wien}!I., Innere Stadt@\textbf{I., Innere Stadt}!Schottenkirche@\textbf{Schottenkirche}, \emph{Kirche}|pw}.\pend
           \vspace{1em}
\pstart
           \raggedleft{}{\pb}10. VIII. 15.\pend
           \vspace{0.5em}
\pstart
           Lieber Arthur, geſtern, an dem erſten leidlichen Tag, habe ich Ihrer
                  Tochter\pwindex{Cappellini, Lili 13.\,9.\,1909 Wien – 26.\,7.\,1928 Venedig@\textsc{Cappellini, Lili} (13.\,9.\,1909 Wien – 26.\,7.\,1928 Venedig)|pwv} einen Beſuch
               gemacht. Sie ſieht wirklich beruhigend aus, hat augenblicklich keine militäriſchen
               Neigungen, will nur junge Dame ſein und hat ſich gnädig von mir die Hand küſſen
               laſſen. Über die \textsc{Ischler Esplanade\oindex{Esplanade [Bad Ischl]@\textbf{Esplanade [Bad Ischl]}, \emph{Straße}|pw}} denkt{ }ſie nicht{ }ſehr gut, hat ſich aber nicht eingehend darüber geäußert, weil
               ihre Aufmerkſamkeit durch \label{K_L04245-1v}\edtext{ihren Hund
               und eine Freundin}{\lemma{\textnormal{\emph{ihren … Freundin}}}\Cendnote{\textnormal{Über einen Hund in
                  ihrem Besitz ist zu dieser Zeit nichts bekannt. Unklar, ob es sich auch bei der
                  Freundin eine fiktive Gefährtin handelt.}}}\label{K_L04245-1} zu ſehr in Anſpruch geno{\geminationm}en war. – –\pend
           
\pstart
           Frau \textsc{Olga\pwindex{Schnitzler, Olga 17.\,1.\,1882 Wien – 13.\,1.\,1970 Lugano@\textsc{Schnitzler, Olga} (17.\,1.\,1882 Wien – 13.\,1.\,1970 Lugano), \emph{Schauspielerin, Sängerin}|pw}}, Sie und \textsc{Heini\pwindex{Schnitzler, Heinrich 9.\,8.\,1902 Hinterbrühl – 12.\,7.\,1982 Wien@\textsc{Schnitzler, Heinrich} (9.\,8.\,1902 Hinterbrühl – 12.\,7.\,1982 Wien), \emph{Regisseur, Schauspieler}|pw}} herzlichſt grüßend\pend
           
\pstart
           Ihr{\\[\baselineskip]}\spacefill\mbox{Gustav.}\pend
           \leftskip=0em{}\selectlanguage{ngerman}\endnumbering\briefempfaengerindex{Schnitzler, Arthur@\textsc{Schnitzler, Arthur}!zzzSchwarzkopf, Gustav@\emph{von Gustav Schwarzkopf}!1915-08-101@{10. 8. 1915}|)be}\mylabel{L04245h}
\begin{anhang}
\end{anhang}\newcommand{\dateiname}{L04245}\newcommand{\titel}{Gustav Schwarzkopf an Arthur Schnitzler, 10. 8. 1915}\newcommand{\editorInnen}{Herausgegeben von Jahnke, SelmaMüller, Martin Anton}\input{../tex-inputs/latex-leseansicht-abspann}
